\section{The rank-one Euler-system hypothesis for the cubic type-II tensor}
\label{sec:cubic-rankin-parity}

The cubic type-II comparison suggests replacing a rational
Dirichlet twist in Kato's input by an Artin tensor product.
This is a specific route that must meet the hypotheses of the
available Euler-system theorem.  Its rank-one image condition
has the following exact obstruction.

\begin{proposition}[AC-K30: orthogonal parity in the cubic Rankin tensor]
\label{AC-K30}
Let $K=\mathbb Q(\omega)$, let $\eta_K$ be its quadratic
character, and let $\chi$ be the cubic character of
$K(\sqrt[3]{269})/K$.  Let $E$ be a field of characteristic
zero containing its values.  Suppose $V$ is any two-dimensional
$E$-representation of $G_{\mathbb Q}$ with
$\det V=\eta_K$.  Put
\[
 U=\operatorname{Ind}_{G_K}^{G_{\mathbb Q}}\chi,
 \qquad R=V\otimes_E U.
                                                        \tag{RS.1}
\]
Then $R$ has a nondegenerate invariant symmetric bilinear form,
and $\det R=1$.  For every $g\in G_{\mathbb Q}$,
\[
               \dim_E R/(g-1)R\in\{0,2,4\}.           \tag{RS.2}
\]
The same conclusion holds after restriction to any subgroup,
duality, or a Tate twist when restricted to the kernel of the
corresponding cyclotomic character.

Consequently no stable lattice in this representation, or in
such a dual or Tate twist, can have a quotient by $g-1$ which
is free of rank one at an element of
$G_{\mathbb Q(\mu_{\ell^\infty})}$.
This excludes the rank-one image hypothesis of the particular
Rankin--Eisenstein Euler-system application described below.
It does not prove that its desired Bloch--Kato Selmer group
is nonzero, or that another Euler-system or deformation
method cannot give its vanishing.
\end{proposition}
\begin{proof}
\proofstep{1. Determinants and the symmetric tensor pairing.}
Choose complex conjugation $c$ with $c^2=1$.
Since $\chi^c=\chi^{-1}$, on $G_K$ the induced
representation has diagonal characters $\chi,\chi^{-1}$,
and $c$ exchanges its two standard coset vectors.
It therefore has determinant one on $G_K$ and determinant
minus one at $c$.  Thus $\det U=\eta_K$.
For any two-dimensional vector space the wedge product is a
nondegenerate alternating pairing with values in its determinant
line.  The tensor of these pairings for $V$ and $U$ is
nondegenerate and symmetric on $R$, since exchanging its two
arguments introduces two minus signs.  Its value line has
character $\det V\det U=\eta_K^2=1$.
Choose a nonzero constant vector in that line to obtain the
claimed invariant scalar pairing $B$.
The tensor determinant formula gives
$\det R=(\det V)^2(\det U)^2=1$.

\proofstep{2. The parity statement without semisimplicity.}
We prove the elementary fact that if $g$ is an isometry of a
nondegenerate symmetric form over a characteristic-zero field,
and $\det g=1$, then $\operatorname{rank}(g-1)$ is even.
No diagonalizability of $g$ is needed.
For an integer $m$ at least the dimension, the primary
decomposition for the factor $X+1$ gives
\[
 E_-=\ker(g+1)^m,\qquad E_0=\operatorname{im}(g+1)^m,
 \qquad R=E_-\oplus E_0.
                                                        \tag{RS.3}
\]
These subspaces are orthogonal.  Indeed the adjoint of $g$
is $g^{-1}$, so for $u\in E_-$,
\[
 B(u,(g+1)^m v)
       =B((g^{-1}+1)^m u,v)
       =B(g^{-m}(g+1)^m u,v)=0.
\]
The restrictions of $B$ to both subspaces are therefore
nondegenerate.  On $E_0$, $g+1$ is invertible, and the
Cayley transform
\[
                     C=(g-1)(g+1)^{-1}
\]
is skew-adjoint for $B$: substitution of $g^*=g^{-1}$
gives $C^*=-C$.
Thus $B(Cx,y)$ is an alternating bilinear form.
Every alternating form in characteristic different from two
has even rank: if it is nonzero, a pair with nonzero pairing
spans a nondegenerate two-dimensional subspace, and induction
on its orthogonal complement proves the assertion.
It follows that $\operatorname{rank}(g-1)|_{E_0}$ is even.
Also $C^*$ is similar to $C^{\rm t}$, so
$\det(1+C)=\det(1-C)$.
Since $g=(1+C)(1-C)^{-1}$ on $E_0$, this proves
$\det(g|_{E_0})=1$.

On $E_-$, $g$ is minus the identity plus a nilpotent
operator.  Hence $g-1$ is invertible there and
$\det(g|_{E_-})=(-1)^{\dim E_-}$.
The assumption $\det g=1$ now forces $\dim E_-$ to be even.
Adding the two ranks proves the elementary fact.
In dimension four the nullity and the coinvariant dimension
are therefore even, proving (RS.2).

\proofstep{3. Lattices and the exact scope of the obstruction.}
Duality preserves the invariant nondegenerate symmetric
form and determinant-one property.  A Tate twist is trivial
on the cyclotomic kernel.  Restricting to a subgroup preserves
all the displayed identities.
If a stable lattice $T$ had $T/(g-1)T$ free of rank one,
tensoring this quotient with its fraction field would give
a one-dimensional quotient of $R$ by $g-1$, contrary to
(RS.2).  The argument uses rational dimension, so enlarging
the coefficient field or changing the stable lattice cannot
remove this obstruction.
\end{proof}

\paragraph{The exact theorem being tested.}
Kings--Loeffler--Zerbes, \emph{Rankin--Eisenstein classes and
explicit reciprocity laws}, arXiv:1503.02888v6, including its
November 2023 corrigendum, Hypothesis~11.1.2(ii), printed
p.~53, requires precisely such a rank-one quotient on the
cyclotomic kernel.  Remark~11.1.3 there already states the
necessary nontriviality of the product of the two nebentypus
characters.  The Bloch--Kato conclusion in Theorem~11.7.3,
printed p.~60, explicitly retains that hypothesis.
\url{https://arxiv.org/pdf/1503.02888}.
The present proposition is a direct proof for the actual
cubic tensor, not a claim that this general literature
obstruction is new.

In the source normalization the descended two-dimensional
type-II representation has determinant $\eta_K$; the
weight-one CM representation $U$ has the same determinant.
The product of the finite nebentypus characters is consequently
one, and the necessary image condition fails even before
testing ordinarity, local Euler factors, or level hypotheses.
The 2026 refinement by Elie Studnia,
\emph{On obstructions to the Euler system method for
Rankin--Selberg convolutions}, Research in Number Theory
12, article~59, Question~1 and Theorem~A, starts with a
nontrivial product and does not remove this obstruction.
\url{https://doi.org/10.1007/s40993-026-00718-1}.
The conditional cubic extraction in AC-K29 remains available,
but this particular replacement of its type-II vanishing input
is not furnished by those theorems.

\paragraph{Dictionary scope.}
Induction, determinant lines and the tensor pairing are
operations on the integral or rational lisse-representation
side of the already audited SGA~1 finite-\'etale dictionary.
The bilinear-form calculation is a coefficient calculation;
it is not a construction of a polarized arithmetic surface
or a new geometric weight theorem.
